% related_work.tex

\section{Related Work}

\paragraph{Multi-turn jailbreak attacks.}
Multi-turn jailbreak families share a persuasion substrate~\citep{zeng2024pap,hagendorff2026reasoning,li2024mhj} but operationalize it differently: gradual escalation~\citep{russinovich2024crescendo}, commitment-consistency bias~\citep{fitd2024}, emotional manipulation~\citep{deceptivedelight2024}, persona steering~\citep{ren2024actorattack}, context poisoning~\citep{alobaid2026echochamber}. Automated methods~\citep{chao2023pair,mehrotra2024tap} iteratively refine adversarial prompts~\citep{perez2022redteaming}. Taxonomies~\citep{wei2023jailbroken} and benchmarks~\citep{mazeika2024harmbench,chao2024jailbreakbench,jiang2024redqueen,song2026multibreak} standardize evaluation. This shared substrate with diverse surface strategies defines our setting: a detector must generalize across families, not merely succeed on its training distribution.

\paragraph{Multi-turn defenses and production safety.}
No existing defense evaluates cross-attack transfer. Recent conversation-level systems pair fine-tuned encoders with sequence aggregation~\citep{deepcontext2026,tong2025biid,gguard2025,safedream2026}, yet each evaluates on a single attack family; absent public implementations preclude direct comparison (Appendix~\ref{app:bert_proxy}). Representation-level approaches~\citep{lu2025xboundary,zhang2025jbshield,xie2024gradsafe,zou2024circuitbreakers,wu2025safeint,shen2024jailbreakantidote,yousefpour2025repbend,kulkarni2026latentadversarial,kadali2025internallayers,kadali2026trace,hua2026rcs} require white-box access and can be brittle under adaptive attacks~\citep{derya2026revisitingjbshield}. Per-turn classifiers~\citep{inan2023llamaguard,metallamaguard3_2024,zeng2024shieldgemma,han2024wildguard,anthropic2025constitutional} and production systems~\citep{markov2023holistic} cannot model multi-turn trajectories; perturbation-based~\citep{robey2024smoothllm} and decoding-based~\citep{xu2024safedecoding} defenses target single-turn attacks. Our ablation fills this gap: decomposition across thirteen configurations under OOD evaluation identifies domain exposure as the dominant factor behind cross-attack transfer.

\paragraph{Computational persuasion and discourse.}
Discourse and persuasion frameworks~\citep{mann1988rst,prasad2008pdtb,stolcke2000dialogueact,austin1962speechacts,searle1969speechacts} offer interpretable vocabulary for attack strategies. Building on Cialdini's principles~\citep{cialdini2001influence}, PAP~\citep{zeng2024pap} identifies 40 persuasion techniques and \citet{hagendorff2026reasoning} find identifiable strategies in 97\% of successful attacks; propaganda analysis~\citep{dasanmartino2019fine,dasanmartino2020semeval} and persuasive dialogue~\citep{wang2019persuasion} apply similar classification. Our ablation disentangles these roles: the taxonomy enables strategy characterization, but detection operates through domain-adapted embeddings regardless of label correctness.

\paragraph{Domain-adaptive pretraining.}
Continued pretraining on domain text yields consistent gains~\citep{gururangan2020dontstop,lee2020biobert,beltagy2019scibert,howard2018ulmfit}, yet this has not been disentangled for safety-critical generalization. Intermediate-task transfer~\citep{phang2018stilts} and multi-task learning~\citep{caruana1997multitask,liu2019mtdnn} confound domain exposure with supervised signal; our MLM control isolates the contribution. Dataset artifact analysis~\citep{gururangan2018artifacts,niven2019probing} motivates the TF-IDF baseline, adversarial NLP~\citep{jin2020textfooler,li2020bertattack,nasr2025attackermovessecond,andriushchenko2024adaptive} informs the perturbation protocol, and probing~\citep{conneau2018probing,hewitt2019structural} underpins CKA analysis.
